For this question use sky map Map-OP02.

The sky projected on the planetarium dome is as seen by an observer on some
hypothetical ringed planet. The observer is located very near to the equator
of the planet.

Assume that the ring lies in the equatorial plane of the planet and the
distance of the planet from its parent star is more than 15\,au. Note that
except for the 10 moons of this planet no other object from this planetary
system (other planets, their satellites, asteroids and even the parent star)
are visible on the projected sky.

At the start of the question the sky will make 3 full rotations at the rate
of 1 rotation every minute. After that the sky will stop rotating at the
observer's midnight. This sky view will remain on the dome for next 7
minutes.

Solve the following questions once the sky stops rotating.

\ioaafig{2025_OBS-PLAN_OP02-f1.pdf}
% exam; not part of the question paper PDF.

\begin{description}
\item[(OP02.1)] Mark the positions of any 4 (only 4) out of the 10 visible
  moons of the planet with a $\otimes$ sign and label them as ``M'' in the
  sky map. \hfill[8]
\item[(OP02.2)] The figure shown below is as seen from the top of the north
  pole of the planet (grey circle) and the horizontal lines are parallel
  light rays from its parent star.

% FIGURE NEEDED: diagram from page 2 of the 2025 planetarium paper — grey
% circle of radius R_P = 6.00 x 10^4 km labelled P, with parallel light rays
% from the parent star drawn as horizontal lines.

  Using the view of the ring on the dome, draw the inner and outer boundaries
  of the ring on this diagram, reproduced in the Summary Answersheet. Your
  sketch must be to scale. \hfill[12]
\item[(OP02.3)] Use the diagram above to determine the width of the ring,
  $w_{\mathrm{ring}}$ (in km), given that the radius of the planet is
  $6.00 \times 10^4$\,km. \hfill[5]
\end{description}
